\section{About this document}

\subsection{Scope and purpose}

This document presents the proof of concept (PoC) of the virtual fence of the DR-VF-01 collar. It
describes what the firmware does, how it was verified and what results were measured, and it states
just as clearly what the PoC does \emph{not} demonstrate. It is an engineering report on one part of
the system; the complete project, with its architecture, hardware, mechanics and compliance, is in the
\emph{Technical Report} (\texttt{docs/DR-VF-01\_Memoria\_Tecnica.pdf}, written in Spanish), which this
report refers to and does not repeat.

The PoC runs on a commercial development board, not on the product's PCB. Every figure in this report
was measured on that board and cannot be transferred as is to the final PCB. See the limitations section.

\subsection{Audience}

It is written for an engineer who wants to understand the design, reproduce the tests or start from this
firmware on another board, and for anyone evaluating the project.

\subsection{Platform and versions}

\begin{table}[H]
\centering
\small
\begin{tabularx}{\textwidth}{|l X|}
\hline
\tablehead{\textbf{Item} & \textbf{Value}} \\ \hline
Board & \dBoard \\ \hline
Application processor and GNSS & \dSoc{} (LTE-M, NB-IoT and GNSS) \\ \hline
SDK & \dNcs{} with \dZephyr \\ \hline
Modem firmware & \dModemFw \\ \hline
Firmware measured & commit \texttt{\dFirmwareCommit} of the repository \\ \hline
Build & \texttt{firmware/build.sh} (see the getting started section) \\ \hline
\end{tabularx}
\caption{Platform and versions of the PoC.}
\end{table}

\subsection{Revision history}

\begin{table}[H]
\centering
\small
\begin{tabularx}{\textwidth}{|c l X|}
\hline
\tablehead{\textbf{Rev.} & \textbf{Status} & \textbf{Content}} \\ \hline
1.0 & Draft & PoC firmware: 60~m square fence created with the first GNSS fix, and a demo mode started with a button. \\ \hline
1.1 & Planned & Polygon defined by user-entered vertices (the fence of a farm). It will be added as a new section and new table rows; it does not rewrite what is already here. \\ \hline
\end{tabularx}
\caption{Report revisions.}
\end{table}

\subsection{How to read the identifiers}

Requirements prefixed \texttt{F-} and \texttt{NF-} are those of the PRD of the Technical Report, and
\texttt{SR-} are those of its SRD. The \texttt{PoC-} criteria belong to this report: each one says which
requirement it derives from, or that it is a new design criterion.
